\section{Visual Prompting：视觉提示的新范式}

\subsection{什么是 Visual Prompting}

Visual Prompting 的核心理念是：\textbf{与其用文字告诉 AI 找什么，不如直接在图片里标注}。具体的视觉线索（Visual Cues）包括：

\begin{itemize}
    \item \textbf{点标注}（Point Prompt）：在图片上指定一个点
    \item \textbf{框标注}（Bounding Box Prompt）：画一个矩形框
    \item \textbf{涂鸦标注}（Scribble）：自由涂画区域
    \item \textbf{文本标注}：在图片上写文字说明
\end{itemize}

这些视觉线索引导模型完成特定的视觉任务，如 VQA、图像分割、目标检测等。在图像生成领域，用户也可以通过在图片中框选区域并写上文字，来指导视觉生成模型做更精确的编辑。

\begin{importantbox}{Visual Prompting 的本质}
Visual Prompting 是一种将人类的视觉意图直接编码到图像中的交互方式。它弥补了纯文本 Prompt 在空间定位上的不足——某些视觉任务（如"最左边的那个蓝色物体"）用文字描述既冗长又容易歧义，而直接在图上标注则直观明了。
\end{importantbox}

\subsection{Visual Prompting 与 Visual Grounding}

Visual Grounding 是 Visual Prompting 的逆过程：
\begin{itemize}
    \item \textbf{Visual Prompting}：人在图上标注 $\rightarrow$ 引导模型理解
    \item \textbf{Visual Grounding}：模型理解文本 $\rightarrow$ 在图上定位（框出/分割出目标）
\end{itemize}

两者结合形成了完整的"视觉对话"闭环。

\subsection{本章小结}

Visual Prompting 将传统的文本提示扩展到视觉维度，为 CV 任务提供了更自然、更精确的人机交互方式。它是理解 SAM3 与 VLM 协作的基础概念。


